\documentclass[12pt]{article}
\usepackage[a4paper,margin=1.25cm]{geometry}
\usepackage{amsmath,amssymb,mathtools}
\usepackage{graphicx}
\usepackage{booktabs}
\usepackage{tabularx}
\usepackage{tikz}
\usepackage[T1]{fontenc}
\usepackage[utf8]{inputenc}
\usepackage[english]{babel}

\setlength{\parindent}{0pt}
\setlength{\parskip}{0.45em}
\setlength{\abovedisplayskip}{5pt plus 1pt minus 1pt}
\setlength{\belowdisplayskip}{5pt plus 1pt minus 1pt}
\setlength{\abovedisplayshortskip}{3pt plus 1pt minus 1pt}
\setlength{\belowdisplayshortskip}{3pt plus 1pt minus 1pt}

\newcommand{\dd}{\,\mathrm{d}}
\newcommand{\xk}{x_k}
\newcommand{\xkp}{x_{k+1}}

\begin{document}

\begin{center}
{\Large\bfseries A cubed-radical sum below \(247/2017\)}\\[0.35em]
{\large two proofs: telescoping after AM--GM, and Jensen against \(x^{-3/2}\)}
\end{center}

\[
\boxed{\displaystyle
\frac1{(1+\sqrt3)^3}
+\frac1{(\sqrt3+\sqrt5)^3}
+\cdots+
\frac1{(\sqrt{2013}+\sqrt{2015})^3}
<\frac{247}{2017}.}
\]

\textbf{Overview.}
Write
\[
x_k=\sqrt{2k-1},\qquad k=1,\dots,1009,
\]
so the sum is \(\sum_{k=1}^{1007}(\xk+\xkp)^{-3}\) with
\(x_1=1\) and \(x_{1008}=\sqrt{2015}\). The whole problem rests on the
identity
\[
\xkp^2-\xk^2=2
\qquad\Longleftrightarrow\qquad
(\xkp-\xk)(\xkp+\xk)=2,
\tag{1}
\]
which converts a sum of \emph{reciprocals} into a sum of
\emph{differences}, the raw material of a telescope. Concretely
\[
\frac1{\xk+\xkp}=\frac{\xkp-\xk}{2},
\qquad\text{hence}\qquad
\frac1{(\xk+\xkp)^3}=\frac{(\xkp-\xk)^3}{8}.
\tag{2}
\]
From here two routes lead to the bound, and both end at the same place:

\begin{itemize}
\item[(A)] AM--GM on \((\xk+\xkp)^2>4\xk\xkp\) turns each term into
\(\frac18\bigl(\frac1{\xk}-\frac1{\xkp}\bigr)\), which telescopes exactly.
\item[(B)] Jensen for \(t\mapsto t^3\) compares
\(\bigl(\xkp-\xk\bigr)^3\) with \(\int_{2k-1}^{2k+1}x^{-3/2}\dd x\),
which telescopes as an integral.
\end{itemize}

The target constant is engineered: \(\frac{247}{2017}=\frac18\bigl(1-\frac{41}{2017}\bigr)\),
and \(\frac{41}{2017}\) is a deliberately weakened form of
\(\frac1{\sqrt{2017}}\), since \(41=\lfloor\sqrt{2017}\rfloor\).

\section*{1. Notation and the basic identity}

Set \(x_k=\sqrt{2k-1}\). Then \(x_1=1\), \(x_2=\sqrt3\), and the sum in
question is
\[
\Sigma=\sum_{k=1}^{1007}\frac1{(\xk+\xkp)^3},
\qquad
\xk+\xkp=\sqrt{2k-1}+\sqrt{2k+1},
\]
the last term being \(k=1007\), i.e. \((\sqrt{2013}+\sqrt{2015})^{-3}\),
as stated.

Since \(\xkp^2-\xk^2=(2k+1)-(2k-1)=2\), we obtain (1) and hence (2). This
already rationalises the denominators: instead of a sum of reciprocals of
sums of surds, we have a sum of cubes of small differences.

\section*{2. Proof A: AM--GM and an exact telescope}

For \(k\ge1\) the numbers \(\xk\ne\xkp\) are positive, so AM--GM is
strict:
\[
\frac{\xk+\xkp}{2}>\sqrt{\xk\xkp}
\qquad\Longrightarrow\qquad
(\xk+\xkp)^2>4\,\xk\xkp
\qquad\Longrightarrow\qquad
\frac1{(\xk+\xkp)^2}<\frac1{4\xk\xkp}.
\tag{3}
\]
Multiply (3) by \(\frac1{\xk+\xkp}\) and use (2):
\[
\frac1{(\xk+\xkp)^3}
<\frac1{4\xk\xkp(\xk+\xkp)}
=\frac{\xkp-\xk}{8\,\xk\xkp}
=\frac18\left(\frac1{\xk}-\frac1{\xkp}\right).
\tag{4}
\]
The right-hand side is a pure difference: the sum telescopes.

It is convenient to add one further term, which is positive and therefore
harmless, so that the telescope ends at \(x_{1009}=\sqrt{2017}\), the
number appearing in the answer:
\[
\Sigma<\sum_{k=1}^{1008}\frac1{(\xk+\xkp)^3}
<\frac18\sum_{k=1}^{1008}\left(\frac1{\xk}-\frac1{\xkp}\right)
=\frac18\left(\frac1{x_1}-\frac1{x_{1009}}\right)
=\frac18\left(1-\frac1{\sqrt{2017}}\right).
\tag{5}
\]

It remains to weaken \(\frac1{\sqrt{2017}}\) into a rational number.
Since \(41^2=1681<2017\), we have \(\sqrt{2017}>41\) and therefore
\[
\frac1{\sqrt{2017}}=\frac{\sqrt{2017}}{2017}>\frac{41}{2017}.
\]
Subtracting a larger quantity gives a smaller result, so from (5)
\[
\Sigma<\frac18\left(1-\frac{41}{2017}\right)
=\frac18\cdot\frac{1976}{2017}
=\frac{247}{2017},
\]
using \(1976=8\cdot247\). \(\square\)

\section*{3. Proof B: Jensen against the integral of \(x^{-3/2}\)}

The second proof replaces AM--GM by convexity and the telescope by the
additivity of an integral.

Since \(\frac{\dd}{\dd x}\sqrt x=\frac1{2\sqrt x}\),
\[
\xkp-\xk=\int_{2k-1}^{2k+1}\frac{\dd x}{2\sqrt x},
\tag{6}
\]
an integral over an interval of length \(2\). Let
\(g(x)=\frac1{2\sqrt x}\) and apply Jensen's inequality in integral form
to the convex function \(\Phi(t)=t^3\) on \((0,\infty)\):
\[
\Phi\!\left(\frac1{b-a}\int_a^b g\right)\le\frac1{b-a}\int_a^b\Phi(g),
\]
with strict inequality because \(g\) is non-constant and \(\Phi\) is
strictly convex. With \(a=2k-1\), \(b=2k+1\), \(b-a=2\), and using (6),
\[
\left(\frac{\xkp-\xk}{2}\right)^3
<\frac12\int_{2k-1}^{2k+1}\frac{\dd x}{8x^{3/2}}
=\frac1{16}\int_{2k-1}^{2k+1}x^{-3/2}\dd x .
\]
By (2) the left-hand side is exactly \(\frac1{(\xk+\xkp)^3}\), so
\[
\boxed{\displaystyle
\frac1{(\xk+\xkp)^3}<\frac1{16}\int_{2k-1}^{2k+1}x^{-3/2}\dd x .}
\tag{7}
\]
Summing (7) for \(k=1,\dots,1007\), the intervals
\([1,3],[3,5],\dots,[2013,2015]\) tile \([1,2015]\) exactly:
\[
\Sigma<\frac1{16}\int_1^{2015}x^{-3/2}\dd x
=\frac1{16}\Bigl[-2x^{-1/2}\Bigr]_1^{2015}
=\frac18\left(1-\frac1{\sqrt{2015}}\right).
\tag{8}
\]
Finally \(45^2=2025>2015\) gives \(\sqrt{2015}<45\), so
\[
\frac1{\sqrt{2015}}>\frac1{45}=0.0222\ldots>\frac{41}{2017}=0.0203\ldots,
\]
and (8) yields
\[
\Sigma<\frac18\left(1-\frac{41}{2017}\right)=\frac{247}{2017}. \qquad\square
\]

\section*{4. How much room there is}

Both proofs are comfortable. The table below shows by how much.

\begin{center}
\small
\renewcommand{\arraystretch}{1.4}
\begin{tabularx}{\textwidth}{@{}>{\raggedright\arraybackslash}p{0.40\textwidth}XX@{}}
\toprule
Quantity & Exact form & Value \\
\midrule
The sum \(\Sigma\) itself
& --
& \(0.1180514\ldots\) \\
Proof B bound, (8)
& \(\frac18\bigl(1-\tfrac1{\sqrt{2015}}\bigr)\)
& \(0.1222153\ldots\) \\
Proof A bound, (5)
& \(\frac18\bigl(1-\tfrac1{\sqrt{2017}}\bigr)\)
& \(0.1222167\ldots\) \\
Target
& \(\frac{247}{2017}\)
& \(0.1224591\ldots\) \\
\bottomrule
\end{tabularx}
\end{center}

So the analytic bounds give about \(0.12222\), against a target of
\(0.12246\): the final rational weakening costs only \(2.4\cdot10^{-4}\),
while the inequality itself has about \(3.5\%\) of slack. The margin is
thin enough that the crude step, replacing \(\sqrt{2017}\) by \(41\),
must be done in the right direction, but comfortable enough that no
sharper estimate is needed.

Note that the two bounds (5) and (8) are numerically almost identical
even though they were obtained by different means and end at different
square roots. That is not a coincidence: both are the same telescoping
estimate of \(\sum\frac18\bigl(x_k^{-1}-x_{k+1}^{-1}\bigr)\), one
performed discretely and one performed as \(\frac1{16}\int x^{-3/2}\),
and the two differ only by one endpoint's worth of the tail.

\section*{5. What each proof really uses}

\begin{center}
\small
\renewcommand{\arraystretch}{1.45}
\begin{tabularx}{\textwidth}{@{}>{\raggedright\arraybackslash}p{0.22\textwidth}XX@{}}
\toprule
& Proof A & Proof B \\
\midrule
Key inequality
& AM--GM: \((\xk+\xkp)^2>4\xk\xkp\)
& Jensen: \(\bigl(\text{mean of }g\bigr)^3<\text{mean of }g^3\) \\
\addlinespace
Telescoping device
& \(\dfrac{\xkp-\xk}{\xk\xkp}=\dfrac1{\xk}-\dfrac1{\xkp}\)
& additivity of \(\displaystyle\int_1^{2015}\) over the intervals \([2k-1,2k+1]\) \\
\addlinespace
Where the exponent \(3\) enters
& once, as the extra factor \(\frac1{\xk+\xkp}\) beyond the squared AM--GM
& directly, as the convex \(\Phi(t)=t^3\) \\
\addlinespace
Generalises to exponent \(p\)?
& only for \(p=3\) in this clean form
& yes: \(\frac1{(\xk+\xkp)^p}<\frac1{2^{p+1}}\displaystyle\int_{2k-1}^{2k+1}x^{-p/2}\dd x\) \\
\bottomrule
\end{tabularx}
\end{center}

The last row is the real advantage of Proof B. Repeating Section 3
verbatim with \(\Phi(t)=t^p\), \(p>1\), gives for every \(n\)
\[
\sum_{k=1}^{n}\frac1{\bigl(\sqrt{2k-1}+\sqrt{2k+1}\bigr)^{p}}
<\frac1{2^{p+1}}\int_1^{2n+1}x^{-p/2}\dd x
=\frac1{2^{p+1}}\cdot\frac{2}{p-2}\left(1-(2n+1)^{1-p/2}\right)
\]
for \(p\ne2\), so the series converges precisely for \(p>2\); at
\(p=3\) the bound is \(\frac18\bigl(1-(2n+1)^{-1/2}\bigr)\), recovering
(8). The exponent \(3\) in the problem is thus the smallest odd integer
for which the sum converges, and the closed form \(\frac18\) of the
limiting bound is what makes a clean rational target possible.

\end{document}
